% Options for packages loaded elsewhere
% Options for packages loaded elsewhere
\PassOptionsToPackage{unicode}{hyperref}
\PassOptionsToPackage{hyphens}{url}
%
\documentclass[
  11pt,
  ignorenonframetext,
  aspectratio=169,
]{beamer}
\newif\ifbibliography
\usepackage{pgfpages}
\setbeamertemplate{caption}[numbered]
\setbeamertemplate{caption label separator}{: }
\setbeamercolor{caption name}{fg=normal text.fg}
\beamertemplatenavigationsymbolsempty
% remove section numbering
\setbeamertemplate{part page}{
  \centering
  \begin{beamercolorbox}[sep=16pt,center]{part title}
    \usebeamerfont{part title}\insertpart\par
  \end{beamercolorbox}
}
\setbeamertemplate{section page}{
  \centering
  \begin{beamercolorbox}[sep=12pt,center]{section title}
    \usebeamerfont{section title}\insertsection\par
  \end{beamercolorbox}
}
\setbeamertemplate{subsection page}{
  \centering
  \begin{beamercolorbox}[sep=8pt,center]{subsection title}
    \usebeamerfont{subsection title}\insertsubsection\par
  \end{beamercolorbox}
}
% Prevent slide breaks in the middle of a paragraph
\widowpenalties 1 10000
\raggedbottom
\AtBeginPart{
  \frame{\partpage}
}
\AtBeginSection{
  \ifbibliography
  \else
    \frame{\sectionpage}
  \fi
}
\AtBeginSubsection{
  \frame{\subsectionpage}
}
\usepackage{iftex}
\ifPDFTeX
  \usepackage[T1]{fontenc}
  \usepackage[utf8]{inputenc}
  \usepackage{textcomp} % provide euro and other symbols
\else % if luatex or xetex
  \usepackage{unicode-math} % this also loads fontspec
  \defaultfontfeatures{Scale=MatchLowercase}
  \defaultfontfeatures[\rmfamily]{Ligatures=TeX,Scale=1}
\fi
\usepackage{lmodern}

\ifPDFTeX\else
  % xetex/luatex font selection
\fi
% Use upquote if available, for straight quotes in verbatim environments
\IfFileExists{upquote.sty}{\usepackage{upquote}}{}
\IfFileExists{microtype.sty}{% use microtype if available
  \usepackage[]{microtype}
  \UseMicrotypeSet[protrusion]{basicmath} % disable protrusion for tt fonts
}{}
\makeatletter
\@ifundefined{KOMAClassName}{% if non-KOMA class
  \IfFileExists{parskip.sty}{%
    \usepackage{parskip}
  }{% else
    \setlength{\parindent}{0pt}
    \setlength{\parskip}{6pt plus 2pt minus 1pt}}
}{% if KOMA class
  \KOMAoptions{parskip=half}}
\makeatother


\usepackage{longtable,booktabs,array}
\newcounter{none} % for unnumbered tables
\usepackage{calc} % for calculating minipage widths
\usepackage{caption}
% Make caption package work with longtable
\makeatletter
\def\fnum@table{\tablename~\thetable}
\makeatother
\usepackage{graphicx}
\makeatletter
\newsavebox\pandoc@box
\newcommand*\pandocbounded[1]{% scales image to fit in text height/width
  \sbox\pandoc@box{#1}%
  \Gscale@div\@tempa{\textheight}{\dimexpr\ht\pandoc@box+\dp\pandoc@box\relax}%
  \Gscale@div\@tempb{\linewidth}{\wd\pandoc@box}%
  \ifdim\@tempb\p@<\@tempa\p@\let\@tempa\@tempb\fi% select the smaller of both
  \ifdim\@tempa\p@<\p@\scalebox{\@tempa}{\usebox\pandoc@box}%
  \else\usebox{\pandoc@box}%
  \fi%
}
% Set default figure placement to htbp
\def\fps@figure{htbp}
\makeatother





\setlength{\emergencystretch}{3em} % prevent overfull lines

\providecommand{\tightlist}{%
  \setlength{\itemsep}{0pt}\setlength{\parskip}{0pt}}



 


\setbeamertemplate{navigation symbols}{}
\setbeamertemplate{footline}[frame number]
\setbeamertemplate{itemize items}[circle]
\setbeamertemplate{itemize subitem}[triangle]
\setbeamercolor{frametitle}{fg=black}
\setbeamerfont{frametitle}{series=\bfseries}
\newcommand{\indep}{\perp\!\!\!\perp}
\usepackage{booktabs}
\usepackage{tikz}
\usetikzlibrary{arrows.meta,positioning,calc}
\tikzset{
  dnode/.style={circle,draw,thick,minimum size=9mm,inner sep=1pt,font=\small},
  unode/.style={circle,draw,thick,dashed,minimum size=9mm,inner sep=1pt,font=\small},
  qnode/.style={circle,draw,thick,dotted,minimum size=9mm,inner sep=1pt,font=\small},
  de/.style={-{Stealth[length=2.2mm]},thick},
  dde/.style={-{Stealth[length=2.2mm]},thick,dashed}
}
\usepackage{fancyvrb}
\fvset{fontsize=\scriptsize}
\makeatletter
\def\verbatim@font{\scriptsize\ttfamily}
\makeatother
\makeatletter
\@ifpackageloaded{caption}{}{\usepackage{caption}}
\AtBeginDocument{%
\ifdefined\contentsname
  \renewcommand*\contentsname{Table of contents}
\else
  \newcommand\contentsname{Table of contents}
\fi
\ifdefined\listfigurename
  \renewcommand*\listfigurename{List of Figures}
\else
  \newcommand\listfigurename{List of Figures}
\fi
\ifdefined\listtablename
  \renewcommand*\listtablename{List of Tables}
\else
  \newcommand\listtablename{List of Tables}
\fi
\ifdefined\figurename
  \renewcommand*\figurename{Figure}
\else
  \newcommand\figurename{Figure}
\fi
\ifdefined\tablename
  \renewcommand*\tablename{Table}
\else
  \newcommand\tablename{Table}
\fi
}
\@ifpackageloaded{float}{}{\usepackage{float}}
\floatstyle{ruled}
\@ifundefined{c@chapter}{\newfloat{codelisting}{h}{lop}}{\newfloat{codelisting}{h}{lop}[chapter]}
\floatname{codelisting}{Listing}
\newcommand*\listoflistings{\listof{codelisting}{List of Listings}}
\makeatother
\makeatletter
\makeatother
\makeatletter
\@ifpackageloaded{caption}{}{\usepackage{caption}}
\@ifpackageloaded{subcaption}{}{\usepackage{subcaption}}
\makeatother

\usepackage{bookmark}
\IfFileExists{xurl.sty}{\usepackage{xurl}}{} % add URL line breaks if available
\urlstyle{same}
\hypersetup{
  pdftitle={Five Common IV Designs},
  pdfauthor={Professor Ji-Woong Chung},
  hidelinks,
  pdfcreator={LaTeX via pandoc}}


\title{\texorpdfstring{Five Common IV Designs}{Five Common IV Designs}}
\subtitle{\texorpdfstring{BUSS975: Causal Inference in Financial
Research}{BUSS975: Causal Inference in Financial Research}}
\author{Professor Ji-Woong Chung}
\date{}
\institute{Korea University}

\begin{document}
\frame{\titlepage}


\begin{frame}{Where we left off}
\protect\phantomsection\label{where-we-left-off}
\begin{itemize}
\tightlist
\item
  We built the machinery: the Wald ratio, 2SLS, weak instruments, and
  the five LATE assumptions. This one is about \textbf{where IVs
  actually come from}.
\item
  IV is a general-purpose strategy: it works wherever a credible shock
  turns up. But over the years certain strategies have been used so
  often that they have become \textbf{designs} --- each with its own
  literature and its own conventions.
\item
  Knowing the designs tells you where to look for an IV in your own
  work.
\end{itemize}
\end{frame}

\begin{frame}{The five}
\protect\phantomsection\label{the-five}
\vspace{0.3em}

\begin{enumerate}
\tightlist
\item
  \textbf{Lotteries} --- randomised offers with imperfect take-up.
\item
  \textbf{Judge / examiner fixed effects} (the ``leniency design'').
\item
  \textbf{Bartik, or shift-share, IVs.}
\item
  \textbf{Fuzzy regression discontinuity} --- which you already know.
\item
  \textbf{Supply shifters for the elasticity of demand} --- the original
  IV, from the 1920s.
\end{enumerate}

\vspace{0.6em}

\begin{itemize}
\tightlist
\item
  For each one, ask the same four questions:
\end{itemize}

\vspace{0.2em}

\begin{center}
\emph{What is the shock? \quad What is the first stage?} 
\end{center}
\begin{center}
\emph{Who are the compliers? \quad Where does exclusion break?}
\end{center}

\vspace{0.3em}

\begin{itemize}
\tightlist
\item
  The last question is the one that decides whether the paper survives
  refereeing.
\end{itemize}
\end{frame}

\section{1. Lotteries}\label{lotteries}

\begin{frame}{The design: randomisation is not enough}
\protect\phantomsection\label{the-design-randomisation-is-not-enough}
\begin{itemize}
\tightlist
\item
  In many randomised trials, assignment is random but \textbf{take-up is
  voluntary}.
\item
  Those who take it up are precisely those who expect to benefit ---
  \textbf{positive selection}.
\item
  So comparing the treated with the untreated is biased \textbf{even
  inside a randomised experiment}.
\end{itemize}

\vspace{0.4em}

\begin{itemize}
\tightlist
\item
  \textbf{The fix:} instrument actual treatment with the
  \textbf{randomised offer}, and estimate the LATE.
\end{itemize}
\end{frame}

\begin{frame}{Oregon Health Insurance Experiment}
\protect\phantomsection\label{oregon-health-insurance-experiment}
\begin{itemize}
\tightlist
\item
  In 2008 Oregon expanded Medicaid for low-income adults. Facing excess
  demand, it allocated places by \textbf{lottery}: roughly \(30{,}000\)
  drawn from a list of \(85{,}000\).
\item
  Of those selected, about \(10{,}000\) enrolled.
\item
  \textbf{First stage} (Finkelstein et al.~2012): winning raised the
  probability of ever being on Medicaid by about \textbf{26 percentage
  points} --- large, but nowhere near one. Hence IV rather than a simple
  comparison.
\item
  \textbf{What the design gives you:} the effect of \emph{having
  Medicaid} for those who enrolled \textbf{because} they won --- not for
  those who would have found coverage anyway.
\end{itemize}
\end{frame}

\begin{frame}{Oregon, in equations}
\protect\phantomsection\label{oregon-in-equations}
\begin{center}\small
\begin{tabular}{ll}
\toprule
$Y_i$ & utilisation or financial strain --- e.g. $\mathbf{1}\{$has a usual place of care$\}$ \\
$D_i$ & $\mathbf{1}\{$enrolled in Medicaid$\}$ \\
$Z_i$ & $\mathbf{1}\{$won the lottery$\}$ \\
\bottomrule
\end{tabular}
\end{center}

\vspace{-0.3em}

\begin{align*}
\text{First stage:}\quad & D_i = \pi_0 + \pi_1 Z_i + \nu_i, \qquad \hat\pi_1 \approx 0.26 \\
\text{Reduced form:}\quad & Y_i = \gamma_0 + \gamma_1 Z_i + u_i \\
\text{Second stage:}\quad & Y_i = \alpha + \delta \hat D_i + e_i, \qquad \hat\delta = \hat\gamma_1 / \hat\pi_1
\end{align*}

\vspace{-0.3em}

\begin{itemize}
\tightlist
\item
  On the next slide, \(\hat\gamma_1\) is the \textbf{ITT column} and
  \(\hat\delta\) is the \textbf{LATE column}.
\end{itemize}
\end{frame}

\begin{frame}{What Oregon found}
\protect\phantomsection\label{what-oregon-found}
\begin{center}
\small
\begin{tabular}{lrr}
\toprule
Outcome & ITT (reduced form) & LATE (2SLS) \\
\midrule
Has a usual place of care      & $9.9\%$  & $33.9\%$ \\
Has a personal doctor          & $8.1\%$  & $28.0\%$ \\
Got all needed healthcare      & $6.9\%$  & $23.9\%$ \\
Any hospital admission         & $0.5\%$  & $2.1\%$  \\
Had a medical collection       & $-1.6\%$ & $-6.4\%$ \\
Owed on medical collection     & $-\$99$  & $-\$390$ \\
\bottomrule
\end{tabular}
\end{center}

\begin{itemize}
\tightlist
\item
  The \textbf{ITT column is the reduced form} --- the policy effect of
  \emph{running the lottery} --- what a legislator decides on.
\item
  The \textbf{LATE column scales it by the first stage} --- the effect
  of \emph{coverage} on the outcomes.
\end{itemize}
\end{frame}

\section{2. Judge and Examiner Fixed
Effects}\label{judge-and-examiner-fixed-effects}

\begin{frame}{The recipe}
\protect\phantomsection\label{the-recipe}
The design needs three ingredients:

\begin{enumerate}
\tightlist
\item
  A \textbf{narrow pipeline} every unit must pass through;
\item
  \textbf{Randomly assigned decision-makers} staffing it;
\item
  \textbf{Discretion} among those decision-makers.
\end{enumerate}

\vspace{0.4em}

\begin{itemize}
\tightlist
\item
  The canonical setting is American courts, where defendants are
  randomly assigned to judges.
\item
  \textbf{Gaudet, Harris and John (1933)} found one judge imprisoned in
  34\% of his cases and another in nearly 60\%, with cases randomly
  rotated between them.
\item
  Their conclusion: a sentence may indicate \emph{``either the
  seriousness of his crime or the severity of the judge.''}
\end{itemize}
\end{frame}

\begin{frame}{One modern example, and how far it travels}
\protect\phantomsection\label{one-modern-example-and-how-far-it-travels}
\begin{itemize}
\tightlist
\item
  \textbf{Stevenson (2018)}, Philadelphia bail magistrates. Randomly
  assigned judges differ in how affordable they set bail, so some
  defendants are detained pre-trial \textbf{purely by draw}.
\item
  She finds pre-trial detention raises conviction by about
  \textbf{13\%}, mostly through guilty pleas, and raises sentence length
  and court fees.
\end{itemize}

\vspace{0.4em}

\begin{itemize}
\tightlist
\item
  The design has spread far beyond courts. Anywhere a queue is staffed
  by rotating officials with discretion:

  \begin{itemize}
  \tightlist
  \item
    patent examiners, disability examiners, loan officers, auditors,
    bankruptcy judges, bank supervisors.
  \end{itemize}
\end{itemize}
\end{frame}

\begin{frame}{Leniency, in equations}
\protect\phantomsection\label{leniency-in-equations}
\begin{center}\small
\begin{tabular}{ll}
\toprule
$Y_i$ & convicted $\{0,1\}$ --- also sentence length, court fees \\
$D_i$ & $\mathbf{1}\{$detained pre-trial$\}$ \\
$Z_i$ & leave-one-out detention rate of the judge who drew case $i$ \\
\bottomrule
\end{tabular}
\end{center}

\vspace{-0.3em}

\begin{align*}
\text{First stage:}\quad & D_i = \pi_1 Z_i + \mathbf{c}_{ct}'\gamma_1 + \nu_i \\
\text{Second stage:}\quad & Y_i = \delta \hat D_i + \mathbf{c}_{ct}'\gamma + e_i
\end{align*}

\vspace{-0.3em}

\begin{itemize}
\tightlist
\item
  \(\mathbf{c}_{ct}\) are \textbf{court \(\times\) period fixed effects}
  --- the cell inside which assignment is random. They are not controls;
  they define the experiment.
\item
  Stevenson's \(\hat\delta\): detention raises conviction by about
  \(13\%\).
\end{itemize}
\end{frame}

\begin{frame}{Building the instrument}
\protect\phantomsection\label{building-the-instrument}
\textbf{Not the judge --- the judge's \emph{leniency}:} \emph{of all the
other cases my judge handled, what fraction did they treat?}

\[
Z_i \;=\; \frac{1}{n_{j(i)} - 1} \sum_{\substack{k \,\neq\, i \\ j(k) \,=\, j(i)}} D_k
\qquad \text{(leave-one-out, or \emph{jackknife}, leniency)}
\]

\vspace{-0.5em}

\begin{center}\footnotesize
\begin{tabular}{ll}
$i$ & the case you are building the instrument \emph{for} \\
$j(\cdot)$ & the judge who drew case $i$ \\
$n_{j(i)}$ & that judge's caseload; subtract $1$ because we dropped case $i$ \\
$D_k$ & whether case $k$ was treated (detained, liquidated, granted) \\
\end{tabular}
\end{center}

\vspace{-0.2em}

\begin{itemize}
\tightlist
\item
  A judge who detains \(70\%\) of \emph{everyone else} is a judge your
  case was unlucky to draw --- and that luck is the only thing the
  instrument uses.
\end{itemize}
\end{frame}

\begin{frame}{Two choices that make it work}
\protect\phantomsection\label{two-choices-that-make-it-work}
\begin{itemize}
\tightlist
\item
  \textbf{Why leave your own case out?} Keeping it in puts \(D_i\) on
  \emph{both} sides: \[
  Z_i^{\text{keep}} \;=\; \underbrace{\frac{1}{n_{j(i)}} \sum_{\substack{k \,\neq\, i \\ j(k) \,=\, j(i)}} D_k}_{\text{the judge}} \;+\; \underbrace{\frac{1}{n_{j(i)}}\, D_i}_{\text{your own case}}
  \]

  \begin{itemize}
  \tightlist
  \item
    \(Z_i^{\text{keep}}\) then inherits a slice of \(D_i\), and so a
    slice of the structural error \(\varepsilon_i\) (how serious the
    case is). \textbf{The instrument is no longer as good as random.}
  \item
    The first stage looks strong for a reason that has nothing to do
    with the judge.
  \item
    The slice is \(1/n_{j(i)}\), so the damage grows when caseloads are
    short \emph{or} when genuine leniency differences are small.
  \end{itemize}
\item
  \textbf{Why residualise first?}

  \begin{itemize}
  \tightlist
  \item
    Assignment is random \emph{within} a court \(\times\) period cell,
    not nationally.
  \item
    So build leniency inside the group, and carry those fixed effects
    into both stages.
  \end{itemize}
\end{itemize}
\end{frame}

\begin{frame}{Where it breaks}
\protect\phantomsection\label{where-it-breaks}
\begin{itemize}
\tightlist
\item
  \textbf{Independence} --- usually fine, assignment really is random.
\item
  \textbf{Exclusion} --- the usual worry. A defendant who draws a judge
  known to be severe expects a worse outcome at trial, so pleads guilty
  --- \emph{whether or not he was detained}. That path runs from \(Z\)
  straight to \(Y\), bypassing \textbf{detention}, which is the
  treatment.
\item
  \textbf{Monotonicity} --- the hardest. Lenient overall, strict on drug
  charges, and the assumption is gone. Imbens and Angrist flagged this
  case in 1994: no nonparametric LATE.
\end{itemize}
\end{frame}

\begin{frame}{Why not just use judge dummies?}
\protect\phantomsection\label{why-not-just-use-judge-dummies}
\begin{itemize}
\tightlist
\item
  What is random is \textbf{which judge you drew}, so 2SLS uses
  \textbf{one dummy per judge} as instruments: \[
  D_i \;=\; \sum_{j=1}^{K} \pi_j\,\mathbf{1}\{j(i) = j\} \;+\; \nu_i \;=\; Z_i'\pi + \nu_i
  \]
\item
  Here \(Z_i\) is the vector of \(K\) judge dummies. The first stage
  sets each \(\hat\pi_j\) to judge \(j\)'s detention rate over
  \textbf{all} her cases --- \textbf{yours included}:
  \(\hat D_i = Z_i'\hat\pi = Z_i^{\text{keep}}\), the instrument we just
  ruled out.
\item
  Count your case, and the judge's ``strictness'' moves with
  \textbf{your own outcome}: a serious case gets detained and makes its
  own judge look strict.
\item
  Leaving it out is the fix. Applied to the dummies' first stage, it is
  called \textbf{JIVE}.
\end{itemize}
\end{frame}

\begin{frame}{How big is the bias? Look at the extremes}
\protect\phantomsection\label{how-big-is-the-bias-look-at-the-extremes}
\begin{itemize}
\tightlist
\item
  \textbf{One case per judge.} The judge's rate \emph{is} your own
  outcome: \(\hat D_i = D_i\). 2SLS \textbf{is} OLS, with all of OLS's
  bias.
\item
  \textbf{Thousands of cases per judge.} Your case is a vanishing slice
  of the rate. 2SLS is the clean IV.
\item
  \textbf{In between}, 2SLS sits between the two, pulled \textbf{toward
  OLS} --- harder the fewer cases each judge has: \[
  \text{2SLS bias} \;\approx\; \frac{\text{OLS bias}}{F+1}
  \]
\item
  \(F\) asks whether judges' rates differ more than chance alone would
  make them. With few cases each, chance alone scatters the rates
  widely, so \(F\) is small. Thin caseloads mean a low \(F\) \emph{and}
  a big slice of your own case --- the same fact, twice.
\end{itemize}
\end{frame}

\begin{frame}{JIVE: leave yourself out}
\protect\phantomsection\label{jive-leave-yourself-out}
\begin{itemize}
\tightlist
\item
  \textbf{The fix:} Estimate the first stage \textbf{without case
  \(i\)}, then use it to predict case \(i\). Same \(K\) instruments ---
  different fitted value. \[
  \hat D_i^{\text{JIVE}} \;=\; Z_i'\hat\pi_{(-i)} , \qquad \hat\pi_{(-i)} \text{ fitted on everyone but } i .
  \]
\item
  Regress \(Y_i\) on \(\hat D_i^{\text{JIVE}}\). This is \textbf{JIVE}
  (Angrist, Imbens and Krueger 1999).
\item
  Nothing about your own case is used to predict your own case.
\end{itemize}

\vspace{0.4em}

\begin{itemize}
\tightlist
\item
  \textbf{You do not run \(n\) regressions.} Leave-one-out has a closed
  form, built from the \textbf{leverage} of observation \(i\): \[
  \hat D_i^{\text{JIVE}} \;=\; \frac{Z_i'\hat\pi - h_i D_i}{1 - h_i},
  \qquad
  h_i = Z_i'(Z'Z)^{-1} Z_i .
  \]
\item
  \(h_i\) is the \(i\)th diagonal of the projection (``hat'') matrix ---
  how much case \(i\) pulls its own fitted value. See Hansen,
  \emph{Econometrics}, §3.21.
\end{itemize}
\end{frame}

\begin{frame}{Your leniency measure \emph{was} JIVE}
\protect\phantomsection\label{your-leniency-measure-was-jive}
With judge dummies and nothing else, your fitted value is just your
judge's average, \(Z_i'\hat\pi = \bar D_{j(i)}\), and your leverage is
your share of that average, \(h_i = 1/n_{j(i)}\), where \(n_{j(i)}\) is
your judge's caseload. So

\[
\hat D_i^{\text{JIVE}}
= \frac{\bar D_{j(i)} - D_i / n_{j(i)}}{1 - 1/n_{j(i)}}
= \frac{1}{n_{j(i)} - 1} \sum_{\substack{k \,\neq\, i \\ j(k) \,=\, j(i)}} D_k .
\]

\begin{itemize}
\tightlist
\item
  That is \textbf{exactly} the leave-one-out leniency we built earlier.
\item
  \textbf{Caveat:} Add controls and this breaks --- next slide.
\end{itemize}
\end{frame}

\begin{frame}{Once you add controls, JIVE breaks too}
\protect\phantomsection\label{once-you-add-controls-jive-breaks-too}
\begin{itemize}
\tightlist
\item
  Assignment is random only \emph{within} a court \(\times\) period
  group, so we control for group fixed effects. By FWL, the instrument
  then becomes \textbf{your judge's rate \(-\) your group's rate}: how
  strict your judge is \emph{relative to the other judges in your
  court}.
\item
  JIVE leaves your case out of the first term only. The group's rate is
  still computed from every case in the group, yours included.
\item
  So your own case walks back in through the second term. The bias
  returns.
\end{itemize}

\vspace{0.4em}

\begin{itemize}
\tightlist
\item
  \textbf{UJIVE} (Kolesár 2013) leaves case \(i\) out of \emph{both}
  pieces: your judge's leave-out rate \(-\) your group's leave-out rate.
\item
  Without group controls the difference is negligible. With them, UJIVE
  is the safe default.
\end{itemize}
\end{frame}

\begin{frame}{Two more things worth knowing}
\protect\phantomsection\label{two-more-things-worth-knowing}
\begin{itemize}
\tightlist
\item
  \textbf{Do not cluster by judge.} It is common, and it is wrong. What
  matters is how cases were \emph{assigned}.

  \begin{itemize}
  \tightlist
  \item
    Cases assigned to judges one at a time, independently
    \(\Rightarrow\) plain \textbf{robust} standard errors.
  \item
    Cases assigned in batches --- one duty judge takes every case filed
    that day \(\Rightarrow\) cluster \textbf{by day}, the unit that was
    actually drawn.
  \end{itemize}
\item
  \textbf{With UJIVE you need not chase a large \(F\).} It stays roughly
  unbiased even when \(F\) is close to \(1\). That is the point of using
  it.
\end{itemize}

\vspace{0.4em}

\begin{itemize}
\tightlist
\item
  \textbf{Read:} Goldsmith-Pinkham, Hull and Kolesár (2026),
  \emph{Leniency Designs: An Operator's Manual}, \emph{JEP} 40(3). A
  step-by-step guide --- including a weaker, more believable version of
  monotonicity and a test for it --- with an \texttt{R} package
  (\texttt{ManyIV}).
\end{itemize}
\end{frame}

\begin{frame}{Stevenson's numbers (Remix, Table 7.12)}
\protect\phantomsection\label{stevensons-numbers-remix-table-7.12}
\begin{center}\small
\begin{tabular}{lcccccc}
\toprule
 & OLS & OLS & 2SLS & 2SLS & JIVE & JIVE \\
\midrule
Detention & $-0.001$ & $0.029$ & $0.151$ & $0.186$ & $0.162$ & $0.212$ \\
          & $(0.002)$ & $(0.002)$ & $(0.065)$ & $(0.064)$ & $(0.070)$ & $(0.076)$ \\
\bottomrule
\end{tabular}
\end{center}

\vspace{-0.2em}
\footnotesize

\(N = 331{,}971\), eight bail judges. Outcome: guilty plea. Odd columns
use time controls only; even columns add defendant and offence
characteristics. \normalsize

\begin{itemize}
\tightlist
\item
  \textbf{OLS says detention does nothing.} IV says it raises guilty
  pleas by \(15\)--\(21\) percentage points.
\item
  But Stevenson has only \textbf{eight} judges. To see what happens with
  \emph{many}, we need a bigger pipeline.
\end{itemize}
\end{frame}

\begin{frame}{A bigger pipeline: patent examiners}
\protect\phantomsection\label{a-bigger-pipeline-patent-examiners}
Goldsmith-Pinkham, Hull and Kolesár (2026) re-run Farre-Mensa, Hegde and
Ljungqvist (2020): does getting your \textbf{first patent approved} help
a startup innovate again?

\begin{itemize}
\tightlist
\item
  \textbf{Treatment} \(D_i\): the application was \textbf{approved}.
\item
  \textbf{Instrument} \(Z_i\): \textbf{one dummy per examiner} ---
  examiners differ in how freely they grant.
\item
  \textbf{Outcome} \(Y_i\): the startup's later innovation
\item
  \textbf{Controls} \(W_i\): \textbf{art unit} \(\times\) year.

  \begin{itemize}
  \tightlist
  \item
    An \emph{art unit} is the small group of examiners who specialise in
    one technology; an application is assigned quasi-randomly to an
    examiner \textbf{within} it. It is the patent office's
    court-and-period group --- the place where the lottery actually
    happens.
  \end{itemize}
\end{itemize}
\end{frame}

\begin{frame}{The standard errors lie too (GPHK 2026)}
\protect\phantomsection\label{the-standard-errors-lie-too-gphk-2026}
\begin{center}\small
\begin{tabular}{lccc}
\toprule
Effect of approval on & UJIVE & 2SLS (examiner dummies) & OLS \\
\midrule
Any later application & $0.173$ & $0.232$ & $0.234$ \\
                      & $(0.055)$ & $(0.016)$ & $(0.006)$ \\
Any later approval    & $0.259$ & $0.240$ & $0.223$ \\
                      & $(0.050)$ & $(0.014)$ & $(0.005)$ \\
$\log(1+\text{citations})$ & $0.419$ & $0.372$ & $0.339$ \\
                      & $(0.125)$ & $(0.033)$ & $(0.011)$ \\
\bottomrule
\end{tabular}
\end{center}

\begin{itemize}
\tightlist
\item
  \textbf{2SLS has landed on OLS} --- That is the bias reproduced.
\item
  \textbf{And look at the standard errors.} 2SLS reports errors three to
  four times \emph{smaller} than UJIVE.
\end{itemize}
\end{frame}

\section{3. Bartik / Shift-Share IVs}\label{bartik-shift-share-ivs}

\begin{frame}{The question: what does immigration do to local wages?}
\protect\phantomsection\label{the-question-what-does-immigration-do-to-local-wages}
\begin{itemize}
\tightlist
\item
  \textbf{The question.} When immigrants arrive in a city, what happens
  to the wages of the \textbf{native} workers already there? \[
  Y_{l,t} = \alpha + \delta D_{l,t} + \rho X_{l,t} + \varepsilon_{l,t}
  \]
\end{itemize}

\vspace{-0.2em}

\begin{itemize}
\tightlist
\item
  \(Y_{l,t}\): log native wages in city \(l\) in year \(t\).
  \(D_{l,t}\): immigrant inflow. \(X_{l,t}\): controls, including city
  and year fixed effects.
\item
  \textbf{Why OLS cannot answer it.} Immigrants are not sprinkled over
  cities at random. They go where the work is.
\item
  \textbf{What we need:} something that sends immigrants to a city for
  reasons that have nothing to do with \emph{that city's} labour market.
\end{itemize}
\end{frame}

\begin{frame}{The idea: you go where your people already are}
\protect\phantomsection\label{the-idea-you-go-where-your-people-already-are}
\begin{itemize}
\tightlist
\item
  A new arrival from Mexico does not choose a city off a map. They go
  where there is already a Mexican community.
\item
  So a \textbf{national} surge in migration from one country lands
  disproportionately in the cities that already held people from that
  country.
\item
  \(l\) = a \textbf{city}, \(k\) = a \textbf{source country}, \(t\) =
  the year, \(t_0\) = a base year \emph{before} the flows we study. \[
  \underbrace{B_{l,t}}_{\text{predicted inflow}} \;=\; \sum_{k=1}^{K} \underbrace{z_{l,k,t_0}}_{\substack{\text{\textbf{shares}} \\ \text{who was already there}}} \times \underbrace{m_{k,t}}_{\substack{\text{\textbf{shifts}} \\ \text{who left home}}}
  \]
\end{itemize}

\vspace{0.2em}

\begin{itemize}
\tightlist
\item
  Los Angeles chose neither the surge out of Mexico nor, in anticipation
  of it, its 1980 immigrants. The surge was decided in Mexico.
\end{itemize}
\end{frame}

\begin{frame}{Shift-share, in equations}
\protect\phantomsection\label{shift-share-in-equations}
\begin{center}\small
\begin{tabular}{ll}
\toprule
$Y_{l,t}$ & log wages of \emph{native} workers in city $l$ \\
$D_{l,t}$ & immigrant inflow into city $l$ --- the Remix writes this $I_{l,t}$ \\
$Z_{l,t}$ & $B_{l,t} = \sum_k z_{l,k,t_0}\, m_{k,t}$ --- inflow predicted by the enclave map \\
$X_{l,t}$ & controls: city and year fixed effects \\
\bottomrule
\end{tabular}
\end{center}

\vspace{-0.3em}

\begin{align*}
\text{First stage:}\quad & D_{l,t} = \pi_0 + \pi_1 B_{l,t} + \rho_1 X_{l,t} + \nu_{l,t} \\
\text{Second stage:}\quad & Y_{l,t} = \alpha + \delta \hat D_{l,t} + \rho X_{l,t} + e_{l,t}
\end{align*}

\vspace{-0.3em}

\begin{itemize}
\tightlist
\item
  First stage: do the cities whose enclaves \emph{should} have drawn
  immigrants actually get them?
\item
  \(\delta\) is the effect of immigration on native wages --- an average
  across source countries, weighted towards whichever ones actually move
  the instrument.
\end{itemize}
\end{frame}

\begin{frame}{Identification}
\protect\phantomsection\label{identification}
\textbf{Bartik IV requires exogeneity of either shares OR shifts}.

\begin{columns}[T]
\begin{column}{0.5\linewidth}
\begin{block}{Framework A: Shares Exogenous}
\protect\phantomsection\label{framework-a-shares-exogenous}
\textbf{Goldsmith-Pinkham, Sorkin, Swift (GPSS, 2020)}

\begin{itemize}
\tightlist
\item
  Identifies via baseline shares \(z_{ik}\).
\item
  Equivalent to using shares as individual instruments.
\item
  Shifts \(m_{k,t}\) act merely as a weighting matrix.
\end{itemize}
\end{block}
\end{column}

\begin{column}{0.5\linewidth}
\begin{block}{Framework B: Shifts Exogenous}
\protect\phantomsection\label{framework-b-shifts-exogenous}
\textbf{Borusyak, Hull, Jaravel (BHJ, 2022)}

\begin{itemize}
\tightlist
\item
  Identifies via sectoral shocks \(m_{kt}\).
\item
  Shares \(z_{ik}\) can be completely endogenous.
\item
  Requires many independent shocks (\(K \to \infty\)).
\end{itemize}
\end{block}
\end{column}
\end{columns}
\end{frame}

\begin{frame}{How identification fails: Framework A (exogenous shares)}
\protect\phantomsection\label{how-identification-fails-framework-a-exogenous-shares}
\begin{itemize}
\tightlist
\item
  \textbf{Path dependence}

  \begin{itemize}
  \tightlist
  \item
    Instrument: 1990 manufacturing shares → local labor demand.
  \item
    Problem: Rust Belt cities with high manufacturing shares were
    already declining, so the shares correlate with \(\varepsilon\)
  \end{itemize}
\end{itemize}

\vspace{0.3em}

\begin{itemize}
\tightlist
\item
  \textbf{Agglomeration and amenities}

  \begin{itemize}
  \tightlist
  \item
    Instrument: tech employment share → local labor demand.
  \item
    Problem: tech-heavy cities have educated workers, costly housing and
    VC that move \(Y\) directly, so the shares correlate with
    \(\varepsilon\)
  \end{itemize}
\end{itemize}
\end{frame}

\begin{frame}{How identification fails: Framework B (exogenous shocks)}
\protect\phantomsection\label{how-identification-fails-framework-b-exogenous-shocks}
\begin{itemize}
\tightlist
\item
  \textbf{Common macro shocks}

  \begin{itemize}
  \tightlist
  \item
    National auto growth spikes because interest rates were cut.
  \item
    Problem: a rate cut raises national auto growth (the shift), and
    auto-heavy cities' spending also responds to the cut directly, so
    the shift correlates with those cities' \(\varepsilon\)
  \end{itemize}
\end{itemize}

\vspace{0.3em}

\begin{itemize}
\tightlist
\item
  \textbf{Few dominant shocks}

  \begin{itemize}
  \tightlist
  \item
    2 of 50 industries (e.g., oil \& gas, healthcare) drive most of the
    shock variation.
  \item
    Problem: BHJ needs many independent shocks so unobservables average
    out. With one or two dominant sectors, any confounder in those
    sectors biases the estimate.
  \end{itemize}
\end{itemize}
\end{frame}

\begin{frame}{Choose your framework first}
\protect\phantomsection\label{choose-your-framework-first}
\begin{itemize}
\tightlist
\item
  \textbf{State your framework:} exogenous shares (GPSS) or exogenous
  shocks (BHJ).
\item
  The framework determines the diagnostics:

  \begin{itemize}
  \tightlist
  \item
    GPSS: Rotemberg weights; balance and pre-trends for the top-weighted
    industries.
  \item
    BHJ: shock-level balance; effective number of shocks;
    exposure-robust SEs.
  \end{itemize}
\end{itemize}

\vspace{0.3em}

\begin{itemize}
\tightlist
\item
  \textbf{Read:} Breuer (2022), \emph{Bartik Instruments: An Applied
  Introduction}, \emph{JFR} 7(1). Written for accounting and finance.
\item
  \textbf{Read:} Borusyak, Hull and Jaravel (2025), \emph{A Practical
  Guide to Shift-Share Instruments}, \emph{JEP} 39(1).
\end{itemize}
\end{frame}

\section{4. Fuzzy Regression
Discontinuity}\label{fuzzy-regression-discontinuity}

\begin{frame}{Lecture 4, revisited}
\protect\phantomsection\label{lecture-4-revisited}
\begin{itemize}
\tightlist
\item
  In a \textbf{sharp} design, \(D_i = \mathbf{1}\{X_i \geq c_0\}\)
  exactly. In a \textbf{fuzzy} design, crossing the cutoff raises the
  probability of treatment without forcing it: \[
  \lim_{x \downarrow c_0} \Pr(D_i = 1 \mid X_i = x) \;\neq\; \lim_{x \uparrow c_0} \Pr(D_i = 1 \mid X_i = x).
  \]
\item
  The estimator is a \textbf{Wald ratio} of two discontinuities: \[
  \delta_{\text{FRD}} = \frac{\lim_{x\downarrow c_0} E[Y|X=x] - \lim_{x\uparrow c_0} E[Y|X=x]}
                           {\lim_{x\downarrow c_0} E[D|X=x] - \lim_{x\uparrow c_0} E[D|X=x]} .
  \]
\item
  A ratio of two discontinuities is a Wald estimator. So a fuzzy RD
  \textbf{is} an IV design.
\end{itemize}
\end{frame}

\begin{frame}{Fuzzy RD is an IV design}
\protect\phantomsection\label{fuzzy-rd-is-an-iv-design}
\begin{itemize}
\tightlist
\item
  Fuzzy RD \textbf{is} IV, with the instrument
  \(Z_i = \mathbf{1}\{X_i \geq c_0\}\):

  \begin{itemize}
  \tightlist
  \item
    Numerator = reduced-form jump in \(Y\); denominator = first-stage
    jump in \(D\).
  \item
    Estimate by 2SLS within a bandwidth, instrumenting \(D\) with \(Z\)
    and controlling for a local polynomial in \(X\) on each side.
  \item
    Same IV assumptions, holding locally at \(c_0\): crossing the cutoff
    affects \(Y\) only through \(D\) (exclusion), and no one is pushed
    \emph{out} of treatment by crossing it (monotonicity). Relevance.
    SUTVA.
  \item
    So \(\delta_{\text{FRD}}\) is a LATE for \textbf{compliers at the
    cutoff}: units whose treatment is switched by crossing \(c_0\).
  \end{itemize}
\end{itemize}
\end{frame}

\begin{frame}{Fuzzy RD, in equations}
\protect\phantomsection\label{fuzzy-rd-in-equations}
\begin{center}\small
\begin{tabular}{ll}
\toprule
$Y_i$ & the outcome \\
$D_i$ & treatment actually **received** --- not assignment \\
$Z_i$ & $\mathbf{1}\{X_i \geq c_0\}$ --- crossing the cutoff \\
\bottomrule
\end{tabular}
\end{center}

\vspace{-0.3em}

\begin{align*}
\text{First stage:}\quad & D_i = \pi_0 + \pi_1 Z_i + f(X_i - c_0) + \nu_i \\
\text{Second stage:}\quad & Y_i = \alpha + \delta \hat D_i + g(X_i - c_0) + e_i
\end{align*}

\vspace{-0.3em}

\begin{itemize}
\tightlist
\item
  \(f\) and \(g\) are the running-variable polynomials, fitted
  \textbf{separately on each side} and only inside the bandwidth, with
  Lecture 4's kernel weights.
\item
  \(\delta\) is a LATE \emph{at the cutoff}, for units whose treatment
  flips there --- doubly local.
\end{itemize}
\end{frame}

\section{5. The Elasticity of Demand}\label{the-elasticity-of-demand}

\begin{frame}{The Identification Problem}
\protect\phantomsection\label{the-identification-problem}
\begin{itemize}
\tightlist
\item
  \textbf{The Target:} A demand curve is a causal object---the effect of
  price (\(P\)) on quantity demanded (\(Q\)).
\item
  \textbf{The Obstacle:} Observed \((P, Q)\) pairs are
  \textbf{equilibrium points} generated by simultaneous shifts in both
  supply and demand.
\item
  Wright's Solution:

  \begin{itemize}
  \tightlist
  \item
    \textbf{History:} Instrumental variables were invented for this
    exact problem by \textbf{Philip Wright (1928)}.
  \item
    \textbf{The Logic:} To trace out a fixed demand curve, you need
    exogenous variation that shifts \textbf{supply only}.
  \end{itemize}
\end{itemize}
\end{frame}

\begin{frame}{The question: how price-sensitive are fish buyers?}
\protect\phantomsection\label{the-question-how-price-sensitive-are-fish-buyers}
\begin{itemize}
\tightlist
\item
  \textbf{The question.} What is the \textbf{price elasticity of demand}
  for fish? \textbf{Kathryn Graddy} set out to measure it at the
  \textbf{Fulton Fish Market}. \textbf{She hand-collected the data},
  driving in from New Jersey before dawn each morning.
\item
  To hold the good fixed, she narrowed it to one product: \textbf{fresh
  whiting of roughly the same size.}
\end{itemize}

\vspace{0.3em}

\begin{itemize}
\tightlist
\item
  \textbf{The obstacle is Wright's.} Each trading day hands you a single
  \((P,Q)\) pair, set by supply and demand \emph{at once}. Regressing
  \(\log Q\) on \(\log P\) traces neither curve.
\item
  \textbf{The instrument:} the weather at sea in the days \emph{before}
  the fish reached the market. (Graddy 1995, 2006)
\end{itemize}
\end{frame}

\begin{frame}{Fulton Fish, in equations}
\protect\phantomsection\label{fulton-fish-in-equations}
\begin{center}\small
\begin{tabular}{ll}
\toprule
$Y_t$ & $\log Q_t$ --- log pounds of whiting sold on day $t$ \\
$D_t$ & $\log P_t$ --- log average price that day \\
$Z_t$ & $\texttt{Stormy}_t$ --- bad weather at sea over the previous two days \\
\bottomrule
\end{tabular}
\end{center}

\vspace{-0.3em}

\begin{align*}
\text{First stage:}\quad & \log P_t = \pi_0 + \pi_1 \texttt{Stormy}_t + \mathbf{d}_t'\gamma_1 + \nu_t, \qquad \hat\pi_1 \approx 0.35 \\
\text{Second stage:}\quad & \log Q_t = \alpha + \delta\, \widehat{\log P_t} + \mathbf{d}_t'\gamma + e_t, \qquad \hat\delta \approx -1.12
\end{align*}

\vspace{-0.3em}

\begin{itemize}
\tightlist
\item
  \(\mathbf{d}_t\) are day-of-week dummies. The unit is a \textbf{day},
  not a buyer.
\item
  Both variables are in \textbf{logs}, so \(\delta\) \emph{is} the
  elasticity --- no conversion required.
\end{itemize}
\end{frame}

\begin{frame}[fragile]{Fulton Fish Market (Remix, Table 7.14)}
\protect\phantomsection\label{fulton-fish-market-remix-table-7.14}
\begin{center}\small
\begin{tabular}{lcc}
\toprule
Dependent variable: $\log Q$ & OLS & 2SLS \\
\midrule
$\log(\text{Price})$ & $-0.563$*** & $-1.119$** \\
 & $(0.152)$ & $(0.443)$ \\
 & & $[-2.273,\; -0.312]$ \\
\midrule
\multicolumn{3}{l}{\emph{First stage}} \\
\texttt{Stormy} & & $0.346$*** \\
 & & $(0.072)$ \\
Effective $F$ & & $22.929$ \\
\midrule
$N$ & $111$ & $111$ \\
\bottomrule
\end{tabular}
\end{center}

\vspace{-0.4em}
\footnotesize

\textbf{Graddy (2006)}, 111 days at the Fulton Fish Market. \(Q\) = log
pounds of whiting sold, \(P\) = log average daily price. \texttt{Stormy}
records bad weather at sea over the previous two days. Robust standard
errors in parentheses, Anderson--Rubin 95\% interval in brackets. All
models control for day-of-week dummies. \(^{*}p<0.1\), \(^{**}p<0.05\),
\(^{***}p<0.01\).

\normalsize
\end{frame}

\begin{frame}{Reading the fish}
\protect\phantomsection\label{reading-the-fish}
\begin{itemize}
\tightlist
\item
  Stormy weather raises price by about \textbf{35\%}. That is a
  \textbf{supply} shock, not a demand shock --- storms make fish harder
  to catch, they do not make New Yorkers hungrier.
\item
  \textbf{OLS says the elasticity is \(-0.56\). IV says \(-1.12\)} ---
  twice as elastic.
\end{itemize}

\vspace{0.4em}

\begin{itemize}
\tightlist
\item
  \textbf{Why does OLS understate it?} Because the observed
  price--quantity pairs are \emph{equilibria}. Both curves move, so the
  scatter traces out neither one.
\end{itemize}
\end{frame}

\begin{frame}{Even here, it is only a LATE}
\protect\phantomsection\label{even-here-it-is-only-a-late}
\begin{itemize}
\tightlist
\item
  The IV elasticity of \(-1.12\) says: \textbf{for compliers}, a 10\%
  storm-induced price rise cuts quantity by about 11\%.
\item
  Now suppose a seller wants to use this number to set a
  profit-maximising price. Demand looks elastic, so cut prices and sell
  more --- right?
\item
  \textbf{No.} That elasticity was identified by \emph{storms}. When the
  firm cuts its price deliberately, it is not using a storm.
\item
  Buyers read a storm-driven price rise as \textbf{temporary} and a
  posted price cut as \textbf{permanent}. Different shock, different
  response, different elasticity.
\end{itemize}
\end{frame}

\begin{frame}{What a LATE is, and is not}
\protect\phantomsection\label{what-a-late-is-and-is-not}
\begin{itemize}
\tightlist
\item
  This is the sobering version of LATE:
\end{itemize}

\vspace{0.2em}

\begin{center}
\emph{It does not tell you what you want to know.}\\
\emph{It tells you what you can learn, given your instrument.}
\end{center}

\vspace{0.4em}

\begin{itemize}
\tightlist
\item
  The warning applies to \textbf{every} design --- the lottery, the
  judge, the national shock, the threshold, the weather.
\item
  Before reporting a LATE, ask one question: \emph{would the policy I
  have in mind move the same people, in the same way, as my instrument
  did?}
\item
  If the answer is no, you have a clean estimate of something other than
  what you were asked.
\end{itemize}
\end{frame}

\begin{frame}{What the five have in common}
\protect\phantomsection\label{what-the-five-have-in-common}
\small

\begin{center}
\begin{tabular}{llll}
\toprule
Design & The shock & Typical first stage & Hardest assumption \\
\midrule
Lotteries        & randomised offer        & strong & exclusion (usually fine) \\
Judge FE         & assigned decision-maker & many weak instruments & \textbf{monotonicity} \\
Bartik           & national shock $\times$ local share & strong & exclusion (shares or shifts) \\
Fuzzy RD         & crossing the cutoff     & jump in $\Pr(D=1)$ & exclusion at the cutoff \\
Elasticity       & supply shifter          & price response & exclusion (demand channel) \\
\bottomrule
\end{tabular}
\end{center}

\vspace{0.4em}

\normalsize

\begin{enumerate}
\tightlist
\item
  \textbf{Instruments are physical events.} A lottery draw, a shift
  schedule, a national policy, a storm. If you cannot describe the
  physical event that assigns treatment in the real world, you do not
  have an IV.
\item
  \textbf{Estimands are Local (LATE).} An IV only identifies the
  treatment effect for the \textbf{compliers}---the specific people
  moved by that exact shock.
\end{enumerate}

\begin{itemize}
\tightlist
\item
  Always name the compliers in plain English before interpreting your
  estimate.
\end{itemize}
\end{frame}

\section{IV in Finance}\label{iv-in-finance}

\begin{frame}{Where the shocks are}
\protect\phantomsection\label{where-the-shocks-are}
Finance has fewer clean lotteries than labour economics. The recurring
sources are:

\begin{itemize}
\tightlist
\item
  \textbf{Regulatory and legal changes} that bind some firms and not
  others --- staggered state laws, index or listing rules, supervisory
  thresholds.
\item
  \textbf{Weather, geography and distance} --- historically popular,
  increasingly contested, because geography correlates with almost
  everything.
\item
  \textbf{Shift-share exposure} --- a firm's or region's pre-existing
  portfolio of industries, suppliers or customers, interacted with
  aggregate shocks.
\item
  \textbf{Randomly assigned decision-makers} --- bankruptcy judges, bank
  examiners, auditors, analysts, patent examiners.
\item
  \textbf{Lotteries and quasi-lotteries} --- oversubscribed allocations,
  IPO allocation rules, random audits.
\end{itemize}

\vspace{0.4em}

\begin{itemize}
\tightlist
\item
  In finance the agents are \textbf{sophisticated, well advised and
  highly motivated}, so the shock you think is exogenous is often
  anticipated, lobbied for, or selected into.
\end{itemize}
\end{frame}

\begin{frame}{Three questions to ask of any finance IV}
\protect\phantomsection\label{three-questions-to-ask-of-any-finance-iv}
\begin{enumerate}
\tightlist
\item
  \textbf{Is the shock really a shock to someone?} Name the units that
  were moved and the units that were not, and say why the difference was
  not their choice. If you cannot describe the assignment mechanism in a
  sentence a practitioner would recognise, you do not have one.
\item
  \textbf{Is there exactly one channel?} This is where most finance IVs
  fail. A legal change that alters takeover exposure also alters
  managerial incentives, disclosure, and the composition of
  shareholders. Distance to a city affects financing \emph{and} labour
  markets, customers and information.
\item
  \textbf{Who are the compliers, and do you care about them?} A
  regulation that binds only small, constrained firms identifies an
  effect for small, constrained firms --- which may be the opposite of
  the effect for the firms the policy debate is about.
\end{enumerate}
\end{frame}

\begin{frame}{The most useful habit}
\protect\phantomsection\label{the-most-useful-habit}
\begin{itemize}
\tightlist
\item
  Write the exclusion restriction down as a \textbf{sentence, in words}
  --- not as an equation, and not as a citation.
\item
  Show it to someone who knows the \textbf{institution} and not the
  econometrics.
\end{itemize}

\vspace{0.4em}

\begin{itemize}
\tightlist
\item
  If they can name a second channel in under a minute, you have your
  answer.
\item
  If you cannot write the sentence at all, you did not have an
  instrument --- you had a variable that happened to be correlated with
  your treatment.
\end{itemize}

\vspace{0.4em}

\begin{itemize}
\tightlist
\item
  Read: Jiang, W. (2017). Have Instrumental Variables Brought Us Closer
  to the Truth? \emph{Review of Corporate Finance Studies} 6, 127--140.
\end{itemize}
\end{frame}

\begin{frame}{References}
\protect\phantomsection\label{references}
\tiny

\begin{itemize}
\tightlist
\item
  Angrist, J., K. Graddy and G. Imbens (2000). The Interpretation of IV
  Estimators in Simultaneous Equations Models. \emph{ReStud} 67,
  499--527.
\item
  Angrist, J., G. Imbens and A. Krueger (1999). Jackknife Instrumental
  Variables Estimation. \emph{Journal of Applied Econometrics} 14,
  57--67.
\item
  Baicker, K. et al.~(2013). The Oregon Experiment --- Effects of
  Medicaid on Clinical Outcomes. \emph{NEJM} 368, 1713--1722.
\item
  Bartik, T. (1991). \emph{Who Benefits from State and Local Economic
  Development Policies?} Upjohn Institute.
\item
  Borusyak, K., P. Hull and X. Jaravel (2022). Quasi-Experimental
  Shift-Share Research Designs. \emph{ReStud} 89, 181--213.
\item
  Borusyak, K., P. Hull and X. Jaravel (2025). A Practical Guide to
  Shift-Share Instruments. \emph{JEP} 39(1), 181--204.
\item
  Breuer, M. (2022). Bartik Instruments: An Applied Introduction.
  \emph{Journal of Financial Reporting} 7(1), 49--67. --- \textbf{read
  this before building a shift-share instrument.}
\item
  Cunningham, S. \emph{Causal Inference: The Remix}, §7.7 --- this deck
  follows it closely.
\item
  Farre-Mensa, J., D. Hegde and A. Ljungqvist (2020). What Is a Patent
  Worth? Evidence from the U.S. Patent ``Lottery''. \emph{Journal of
  Finance} 75(2), 639--682.
\item
  Finkelstein, A. et al.~(2012). The Oregon Health Insurance Experiment:
  Evidence from the First Year. \emph{QJE} 127, 1057--1106.
\item
  Gaudet, F., G. Harris and C. St.~John (1933). Individual Differences
  in the Sentencing Tendencies of Judges. \emph{Journal of Criminal Law
  and Criminology} 23, 811--818.
\item
  Goldsmith-Pinkham, P., P. Hull and M. Kolesár (2026). Leniency
  Designs: An Operator's Manual. \emph{JEP} 40(3), 213--240. ---
  \textbf{read this before running a leniency design.}
\item
  Goldsmith-Pinkham, P., I. Sorkin and H. Swift (2020). Bartik
  Instruments: What, When, Why, and How. \emph{AER} 110, 2586--2624.
\item
  Graddy, K. (1995). Testing for Imperfect Competition at the Fulton
  Fish Market. \emph{RAND Journal of Economics} 26(1), 75--92.
\item
  Graddy, K. (2006). Markets: The Fulton Fish Market. \emph{JEP} 20,
  207--220.
\item
  Hansen, B. (2022). \emph{Econometrics}. Princeton University Press.
  (§3.21, leave-one-out and leverage.)
\item
  Imbens, G. and J. Angrist (1994). Identification and Estimation of
  Local Average Treatment Effects. \emph{Econometrica} 62, 467--475.
\item
  Kolesár, M. (2013). Estimation in an Instrumental Variables Model with
  Treatment Effect Heterogeneity. Working paper. (UJIVE.)
\item
  Jiang, W. (2017). Have Instrumental Variables Brought Us Closer to the
  Truth? \emph{Review of Corporate Finance Studies} 6, 127--140.
\item
  Stevenson, M. (2018). Distortion of Justice: How the Inability to Pay
  Bail Affects Case Outcomes. \emph{Journal of Law, Economics, and
  Organization} 34, 511--542.
\end{itemize}
\end{frame}




\end{document}
